\subsection{Preprocessing}
\label{sec:preprocessing}
The EDA shows that the images already share a fixed resolution and class
mapping, but differ in brightness, coverage, and spatial appearance. Our
preprocessing therefore keeps the original image content and applies a common
numerical transformation before arranging the pixels for each model. The
pipeline has four steps: convert the raw image to a floating-point tensor,
scale its intensities, standardize it using training statistics, and construct
the required input layout. The same fitted transformation is used for
training, validation, and test images.

\subsubsection{Tensor Conversion and Intensity Scaling}
Let \(X_i\in\{0,\ldots,255\}^{28\times28}\) be a raw image. The loader
uses \texttt{ToTensor} to add a channel dimension, convert the unsigned
8-bit values to \texttt{float32}, and divide the intensities by 255:
\begin{equation}
 U_{i,0,r,c}=\frac{X_{i,r,c}}{255},\qquad
 U_i\in[0,1]^{1\times28\times28}.
 \label{eq:scaling}
\end{equation}
Here, \(r,c\in\{0,\ldots,27\}\) denote the row and column, and the
channel index is zero because the image is grayscale. For example, raw
intensities 0, 128, and 255 become approximately 0, 0.5020, and 1.
This conversion changes the numeric range while preserving the 784 pixels,
their order, and the supplied label. Native resolution is retained; no
interpolation or resizing is needed. Scaling takes place on the CPU in the
dataset loader, before a batch is transferred to the selected training device.

\subsubsection{Standardization Using Training Statistics}
Scaling to \([0,1]\) and standardization serve different purposes. Scaling
sets a common numerical range; standardization centers that range relative
to the training distribution and expresses variation in units of its standard
deviation. We use one mean and one population standard deviation for the
single grayscale channel, computed across all training images and positions.
With \(\mathcal{I}_{\mathrm{tr}}\) denoting the frozen set of 54,000 training
indices and \(N=54{,}000\times784=42{,}336{,}000\), the statistics are
\begin{equation}
 \mu=\frac{1}{N}\sum_{i\in\mathcal{I}_{\mathrm{tr}}}
       \sum_{r=0}^{27}\sum_{c=0}^{27}U_{i,0,r,c},\qquad
 \sigma=\sqrt{\frac{1}{N}\sum_{i\in\mathcal{I}_{\mathrm{tr}}}
       \sum_{r=0}^{27}\sum_{c=0}^{27}(U_{i,0,r,c}-\mu)^2}.
 \label{eq:normalization_stats}
\end{equation}
The measured values are \(\mu=\TrainMean\) and \(\sigma=\TrainStd\),
rounded here to six decimal places. The implementation retains the full
precision stored in the split cache. Each scaled image is transformed as
\begin{equation}
 Z_{i,0,r,c}=\frac{U_{i,0,r,c}-\mu}{\sigma}.
 \label{eq:normalization}
\end{equation}
These are global training-pixel statistics, rather than separate statistics
for each image, class, or batch. Applying the same affine transformation to
every image preserves differences in brightness between examples. Computing
separate statistics for each class would also require the unknown target
label when preprocessing a new image, so it is not part of this pipeline.

Table~\ref{tab:preprocessing_scales} summarizes the three numerical stages.
After standardization, raw zero-valued background pixels become approximately
\(-0.810\), and raw values of 255 become approximately \(2.021\). Negative
values are expected and are retained. The transformation is invertible up to
floating-point rounding: \(U=\sigma Z+\mu\). It does not make the image
distribution Gaussian or remove the background.

\begin{table}[H]
\centering
\caption{Numerical scales across all training pixels. Values are rounded;
standardization is global, not performed separately for each image.}
\label{tab:preprocessing_scales}
\begin{tabular}{lrrrr}
\toprule
Representation & Mean & Population SD & Minimum & Maximum\\
\midrule
Raw intensity & 72.983297 & 90.055735 & 0 & 255\\
Scaled intensity & 0.286209 & 0.353160 & 0 & 1\\
Standardized intensity & $\approx0$ & $\approx1$ & $-0.810424$ & $2.021156$\\
\bottomrule
\end{tabular}
\end{table}

\subsubsection{Preventing Leakage Between Partitions}
The split described in Section~2.2 is fixed before fitting the normalization.
Only its 54,000 training images contribute to \(\mu\) and \(\sigma\).
The validation and test transforms reuse those constants without recalculation;
their own normalized distributions are not required to have exactly zero mean
or unit standard deviation. Fitting statistics on all 70,000 images, or
normalizing each evaluation partition using its own statistics, would allow
held-out information to influence the input transformation.

The cached training and validation indices, normalization constants, source
fingerprint, and split fingerprint identify the preprocessing used by a run.
All five architectures use this same split and fitted numerical transform.
Validation loss selects the checkpoint, while the official test set evaluates
the selected weights. Neither test labels nor test performance determine the
normalization or the choice of representation in the baseline comparison.

\subsubsection{Input Representations for the Five Models}
After standardization, the representation utility arranges the same 784
values for each architecture. Table~\ref{tab:representations} gives the
shapes after batching, where \(B\) is the number of images in a batch.
The normal batch size is 128; the final batch may contain fewer images.
Inputs remain \texttt{float32}. Each image keeps one integer target in
\(\{0,\ldots,9\}\), giving a target tensor of shape \((B,)\).
Labels are not standardized or converted into image features.

\begin{table}[H]
\centering
\caption{Model input layouts after common standardization. Shapes precede
learned feature extraction or patch embedding. Each image retains 784 values.}
\label{tab:representations}
\begin{tabularx}{\linewidth}{l l X}
\toprule
Model & Batch shape & Ordering and interpretation\\
\midrule
Linear / MLP & \((B,784)\) & One vector per image, flattened row by row.\\
CNN & \((B,1,28,28)\) & One grayscale channel with the original image grid.\\
LSTM & \((B,28,28)\) & 28 row steps, with 28 pixel features at each step.\\
Transformer & \((B,49,16)\) & 49 nonoverlapping $4\times4$ patches,
with 16 pixel features per patch.\\
\bottomrule
\end{tabularx}
\end{table}

\textbf{Linear Classifier and MLP.}
The image is flattened in row-major order:
\begin{equation}
 v_{i,28r+c}=Z_{i,0,r,c},\qquad v_i\in\mathbb{R}^{784}.
 \label{eq:flattening}
\end{equation}
Thus, positions 0--27 contain the first row and positions 28--55 contain
the second. Flattening preserves the values and their fixed positions, but
the model interface does not explicitly encode a two-dimensional neighborhood
or impose local weight sharing. The MLP's hidden transformations belong to
the model architecture, rather than to preprocessing.

\textbf{Convolutional Neural Network.}
The CNN receives \(Z_i\) directly as a \((1,28,28)\) image. Keeping the
channel and spatial axes allows convolutional layers to operate on local
neighborhoods. No handcrafted edge map, threshold mask, or class prototype
is supplied as an additional channel.

\textbf{LSTM.}
The single channel dimension is removed, giving
\begin{equation}
 S_{i,t,j}=Z_{i,0,t,j},\qquad t,j\in\{0,\ldots,27\}.
 \label{eq:rows}
\end{equation}
The sequence moves from the top row to the bottom row, and each step contains
the 28 left-to-right pixel values in that row. These are spatial steps within
one static image. Because both axes have length 28, an accidental transpose
would preserve the shape while changing the sequence; ordering must therefore
be checked by comparing the actual elements.

\textbf{Transformer.}
The image is divided into a \(7\times7\) grid of nonoverlapping
\(4\times4\) patches. Patches are read left to right and then top to bottom,
and the pixels inside each patch are also flattened row by row. If
\(a,b\in\{0,\ldots,6\}\) index the patch grid and
\(u,v\in\{0,\ldots,3\}\) index a pixel within a patch, then
\begin{equation}
 P_{i,7a+b,\,4u+v}=Z_{i,0,\,4a+u,\,4b+v}.
 \label{eq:patches}
\end{equation}
There are 49 patches with 16 values each, so \(49\times16=784\).
The patch layout requires no padding or discarded boundary pixels.
Projection from 16 pixel values to an embedding, positional encoding, and
attention are model operations to be explained in the methodology section.

Figure~\ref{fig:representations} illustrates the layouts using one training
image. Reordering preserves the pixel information; differences in how the
architectures learn from these layouts will be examined experimentally.
\reportfigure{07_preprocessing_representations}{1.0}{A training image shown as
raw pixels and as standardized inputs for the five models. Vector, row, and
patch layouts contain the same 784 standardized values; only their arrangement
changes. The example is not a model prediction.}{fig:representations}
\FloatBarrier

\subsubsection{Data Retention and Augmentation Policy}
The baseline retains every image and all pixel coordinates. The audit found
no entirely blank training images or exact duplicates across the checked
partitions. Brightness and nonzero coverage vary within each class, but
extreme values alone do not establish corruption or an incorrect label.
Accordingly, the EDA does not justify removing dark images, cropping stable
background locations, or relabeling visually ambiguous garments.

No resizing, denoising, binarization, histogram equalization, or dimensionality
reduction is applied to the model inputs. In particular, the nonzero masks,
class means, and similarity matrix in the EDA remain descriptive analyses;
they are not extra training features or replacements for the original images.
All five models therefore start from the same grayscale information.

Random augmentation is disabled in the shared baseline configuration. This
keeps the initial comparison independent of architecture-specific augmentation
choices. The loader supports optional horizontal flips with probability 0.5
and rotations of up to ten degrees for training images, but these options are
not part of the reported preprocessing baseline. If studied later, they must
be identified as separate experiments, applied only to training inputs before
standardization, and evaluated with the same deterministic validation and
test transforms. No augmentation gain is assumed here.

\subsubsection{Preprocessing Verification}
The saved EDA audit checks the transformation against the existing data
utility. Over all standardized training pixels, the measured mean is within
\(10^{-6}\) of zero and the population standard deviation is within
\(10^{-6}\) of one. Independent NumPy calculations reproduce the raw moments
and pixel histogram, and the fitted normalization agrees with the split cache.
The same CPU tensor conversion is used by the loader and the EDA before CUDA
calculations, avoiding a change in conversion order between them.

Representation checks compare flattened and row-sequence elements with their
source pixels, reconstruct the image from the patch sequence, and confirm
agreement between CPU and CUDA layout conversions. These checks verify
values and ordering, rather than relying only on tensor dimensions.
Together with the fixed split, they establish the input interface for the
methodology and implementation sections. They do not establish classifier
accuracy, which requires evaluation of trained models.
\FloatBarrier
